\section{Introduction}

Let $\Omega$ be a set and let $\Lo \subseteq \PO(\Omega)$ be a collection of
subsets containing $\varnothing$ and $\Omega$, closed under complementation
$A \mapsto A\orth := \Omega \setminus A$, and closed under \emph{orthogonal}
(pairwise-disjoint) countable unions: if $\{A_n\} \subseteq \Lo$ are pairwise
disjoint then $\bigcup_n A_n \in \Lo$. Ordered by inclusion, such an $\Lo$ is a
\emph{concrete $\sigma$-complete orthomodular lattice} (concrete $\sigma$-OML), in
the sense of Gudder's concrete quantum logics~\cite{Gudder1969,
PtakPulmannova1991}; it is exactly a \emph{$\sigma$-class} in the terminology of
Navara and Pt\'ak~\cite{NavaraPtak1983}. Such structures model the event
algebras of physical systems whose measurement contexts are mutually
incompatible: the maximal Boolean subalgebras (``blocks'') overlap, but their
union need not be Boolean.

A \emph{two-valued state} on $\Lo$ is a map $s\colon \Lo \to \{0,1\}$ with
$s(\Omega) = 1$ that is finitely additive on orthogonal pairs
($s(A) + s(A\orth) = 1$). It is \emph{$\sigma$-additive} if, for every pairwise
disjoint $\{A_n\} \subseteq \Lo$, one has $s\!\left(\bigcup_n A_n\right) =
\sum_n s(A_n)$. For each point $\omega \in \Omega$ the \emph{point evaluation}
$\delta_\omega(A) := \mathbbm{1}[\omega \in A]$ is a $\sigma$-additive two-valued
state; we call these \emph{Dirac} states and all others \emph{non-Dirac}.

The objects of interest are local patterns that cannot be globalised. Fix a
\emph{finite, $\perp$-closed sub-orthoposet} $B \subseteq \Lo$ (a finite
subfamily closed under complementation) and a two-valued state $s_0$ on $B$. We
say $s_0$ \emph{extends} to a global state $s$ on $\Lo$ if $s\restr B = s_0$.

\begin{definition}[$\sigma$-essential contextual state]\label{def:sigma-essential}
A \emph{$\sigma$-essential contextual state} on $\Lo$ is a two-valued state $s_0$
on a finite $\perp$-closed sub-orthoposet $B \subseteq \Lo$ that extends to
\emph{no} global $\sigma$-additive two-valued state on $\Lo$. A concrete
$\sigma$-OML $\Lo$ \emph{carries a $\sigma$-essential contextual state} (is a
\emph{witness}) if such an $s_0$ exists. We require $\Lo$ to be non-Boolean and
\emph{irreducible} (centre $\{\varnothing, \Omega\}$); see
Remark~\ref{rem:offcenter}.
\end{definition}

The motivating phenomenon is a failure of compactness at the $\sigma$-level. By a
theorem of Wright~\cite{Wright1978}, every two-valued state on a \emph{finite}
OML is itself realised by finitely additive data, so no obstruction can be
witnessed finitely; the question is whether a coherent \emph{finite} pattern can
fail to globalise once countable additivity is imposed. On a Boolean algebra the
answer is no, for two independent reasons recalled in \S\ref{sec:boolean}. The
present paper isolates exactly what would have to differ in the non-distributive
case, reduces the question to its irreducible core, and locates that core against
the literature.

\paragraph{Results.}
\S\ref{sec:boolean} records the Boolean baseline. \S\ref{sec:localization}
proves the \emph{Localization Theorem} (Theorem~\ref{thm:localization}): a witness
exists iff (i) no Dirac extends $s_0$ and (ii) no non-Dirac $\sigma$-additive
two-valued state does, and clause (i) is freely arrangeable, so the entire
difficulty localises to clause (ii). \S\ref{sec:bare} strips clause (ii) to a
bare set-theoretic existence statement (a coherent $\sigma$-additive selection
across overlapping countable partitions). \S\ref{sec:boundary} assembles the
boundary map: the question is empty in the Polish-representable case, does not
reduce to the Hilbert-space dichotomy of Blecher--Weaver, and does not reduce to a
measurable cardinal as its Boolean shadow would. \S\ref{sec:open} states the open
problem.

\begin{remark}[Irreducibility is not vacuous]\label{rem:offcenter}
The irreducibility requirement in Definition~\ref{def:sigma-essential} is
compatible with concreteness and is genuinely restrictive rather than
contradictory. The horizontal sum $MO_\kappa$ (the OML with $\kappa$ orthogonal
atom-pairs, $0$, and $1$) is concrete, non-Boolean, and has centre
$\{\varnothing,\Omega\}$ for $\kappa \geq 2$~\cite{Kalmbach1983}; for
$\kappa = \omega$ it is $\sigma$-complete. So the carrier class of
Definition~\ref{def:sigma-essential} is inhabited. Concreteness does \emph{not}
force a non-trivial centre: by the Zierler--Schlessinger/Gudder criterion a
concrete OML is precisely one with an order-determining family of two-valued
states~\cite{Gudder1969,Harding2004}, a state-separation condition unrelated to
reducibility. We exclude $MO_\kappa$ and other \emph{segregated} carriers as
witnesses for an independent reason (Remark~\ref{rem:segregated}).
\end{remark}


\section{The Boolean baseline}\label{sec:boolean}

When $\Lo$ is a Boolean $\sigma$-algebra of subsets of $\Omega$,
Definition~\ref{def:sigma-essential} is never satisfied, for two reasons that we
isolate because the non-distributive case must defeat both.

\begin{proposition}[No witness in the Boolean case]\label{prop:boolean}
Let $\Lo \subseteq \PO(\Omega)$ be a Boolean $\sigma$-algebra and $B \subseteq
\Lo$ a finite $\perp$-closed sub-orthoposet with a two-valued state $s_0$. Then
$s_0$ extends to a global $\sigma$-additive two-valued state on $\Lo$.
\end{proposition}

\begin{proof}
Since $B$ is finite, the family $\{A \in B : s_0(A) = 1\}$ is finite; because
$s_0$ is a state on $B$ and $B$ is $\perp$-closed, this family has the finite
intersection property and its intersection $K := \bigcap\{A \in B : s_0(A)=1\}$
lies in the Boolean algebra. If $K \neq \varnothing$, any $\omega \in K$ gives
$\delta_\omega \restr B = s_0$ (see Lemma~\ref{lem:dirac-iff} below), and
$\delta_\omega$ is $\sigma$-additive. If $K = \varnothing$, then $B$ exhibits a
finite cover by complements and the assumption that $s_0$ is a coherent state
forces an atom of the finite Boolean subalgebra generated by $B$ on which $s_0$
concentrates; extend by the corresponding point or, in the atomless case, by the
fact that a two-valued $\sigma$-additive measure on a Boolean $\sigma$-algebra
exists on the generated $\sigma$-subalgebra and, absent a measurable cardinal, is
a point mass~\cite{Sikorski1964,Jech2003}. Either way $s_0$ globalises.
\end{proof}

The two reasons are visible in the proof: (a) the local pattern is determined by
\emph{finitely many compatible} events, which on a Boolean algebra always admit a
common point or atom; and (b) the only global $\sigma$-additive two-valued states
on a Boolean $\sigma$-algebra are, away from a measurable cardinal, Dirac. A
non-distributive witness must break (a) --- the events of $B$ live in
\emph{incompatible} blocks with no common point --- while \emph{also} ensuring
(b) cannot be repaired by a non-Dirac global state. The tension between these is
the subject of the paper.


\section{The Localization Theorem}\label{sec:localization}

We first record the elementary characterisation of when a Dirac extends $s_0$.

\begin{lemma}\label{lem:dirac-iff}
Let $B \subseteq \Lo$ be a finite $\perp$-closed sub-orthoposet with two-valued
state $s_0$. A point evaluation $\delta_\omega$ satisfies $\delta_\omega\restr B
= s_0$ iff $\omega \in K(s_0) := \bigcap\{A \in B : s_0(A) = 1\}$. Consequently no
Dirac extends $s_0$ iff $K(s_0) = \varnothing$.
\end{lemma}

\begin{proof}
$\delta_\omega\restr B = s_0$ means $\omega \in A \iff s_0(A)=1$ for every
$A \in B$. The forward constraints place $\omega \in \bigcap\{A : s_0(A)=1\}$. For
a false $A$ (i.e.\ $s_0(A)=0$), $\perp$-closure gives $A\orth \in B$, and
additivity gives $s_0(A\orth) = 1 - s_0(A) = 1$; so $A\orth$ is one of the true
elements and the constraint $\omega \notin A$ is already $\omega \in A\orth$,
subsumed in $K(s_0)$. Thus the false constraints add nothing, and
$\delta_\omega \restr B = s_0 \iff \omega \in K(s_0)$. Conversely any
$\omega \in K(s_0)$ works: true $A$ give $\omega\in A$, false $A$ give
$\omega \in A\orth$ hence $\omega\notin A$.
\end{proof}

\begin{theorem}[Localization]\label{thm:localization}
Let $\Lo$ be a concrete $\sigma$-OML, $B \subseteq \Lo$ a finite $\perp$-closed
sub-orthoposet, and $s_0$ a two-valued state on $B$. Then $s_0$ is a
$\sigma$-essential contextual state iff
\begin{enumerate}[label=\textup{(\roman*)}]
  \item $K(s_0) = \varnothing$ \textup{(}no point evaluation extends
    $s_0$\textup{)}, and
  \item no \emph{non-Dirac} $\sigma$-additive two-valued state on $\Lo$ extends
    $s_0$.
\end{enumerate}
Moreover clause \textup{(i)} is freely arrangeable: there is a concrete
$\sigma$-OML and an $s_0$ with $K(s_0)=\varnothing$.
\end{theorem}

\begin{proof}
Every global $\sigma$-additive two-valued state on $\Lo$ is either a Dirac
$\delta_\omega$ or non-Dirac, and this dichotomy is exhaustive. Hence ``no global
$\sigma$-additive two-valued state extends $s_0$'' is equivalent to the
conjunction ``no Dirac extends $s_0$'' and ``no non-Dirac extends $s_0$.'' By
Lemma~\ref{lem:dirac-iff} the first conjunct is (i); the second is (ii). For
arrangeability of (i), the Navara--Pt\'ak construction
(Example~\ref{ex:navara-ptak}) realises $K(s_0)=\varnothing$ explicitly.
\end{proof}

\begin{remark}[What Theorem~\ref{thm:localization} is, and is not]\label{rem:localization-scope}
Theorem~\ref{thm:localization} is a decomposition, not a reduction to a separate,
harder named problem. Because the Dirac/non-Dirac split is exhaustive, the
equivalence ``$s_0$ extends to no $\sigma$-additive two-valued state iff (i) and
(ii)'' merely unfolds the definition of witness. Its content is the second clause
of the statement: clause (i) --- the obstruction visible to point evaluations ---
is \emph{freely arrangeable}, so the entire mathematical and set-theoretic weight
of Definition~\ref{def:sigma-essential} sits in clause (ii). We therefore call
clause (ii) the \emph{$\sigma$-point-selection core} and devote the remainder of
the paper to it.
\end{remark}

The role of the Navara--Pt\'ak construction is twofold, and it is worth being
precise, because the same example both supplies clause (i) and shows clause (ii)
is the binding one.

\begin{example}[Navara--Pt\'ak~\cite{NavaraPtak1983}]\label{ex:navara-ptak}
Let $\Omega = \bbQ_0 \times \bbQ_0$ with $\bbQ_0 = \bbQ \cap (0,1)$, and let
$f,g\colon \Omega \to \bbR$ be the coordinate projections. Let $\Lo =
\Lo_{f,g}$ be the least $\sigma$-class containing the inverse images of Borel
sets under $f$ and $g$. Set $C_f = f^{-1}(\{1/3\})$, $C_g = g^{-1}(\{1/2\})$,
$C_{f+g} = (f+g)^{-1}(\{1/2\})$. Define $m(A) = 1$ iff $A$ contains one of
$C_f, C_g, C_{f+g}$. Navara and Pt\'ak show $m$ is a (two-valued, $\sigma$-additive)
state on $\Lo$, and that since $C_f \cap C_g \cap C_{f+g} = \varnothing$ the state
$m$ is \emph{not} concentrated at any point.
\end{example}

\begin{remark}[The example is not itself a witness]\label{rem:np-builds-rescuer}
In the notation of Theorem~\ref{thm:localization}, take $B$ to be the
$\perp$-closure of $\{C_f, C_g, C_{f+g}\}$ and $s_0$ the pattern sending each to
$1$. Because $C_f \cap C_g \cap C_{f+g} = \varnothing$, no point evaluation
extends $s_0$: clause (i) holds. But Navara and Pt\'ak's $m$ is a non-Dirac
$\sigma$-additive two-valued state with $m\restr B = s_0$; clause (ii)
\emph{fails}, so $s_0$ is not $\sigma$-essential. The construction \emph{builds
the rescuer}. This is exactly why clause (i) alone never suffices and the whole
question is clause (ii): one must produce a carrier on which the point-level
obstruction \emph{cannot} be repaired by any non-Dirac state. Their carrier
$\Lo_{f,g}$, being generated by two compatible (commuting) observables, is
Boolean enough to admit the repair.
\end{remark}


\section{The bare statement}\label{sec:bare}

We now remove all lattice scaffolding from the $\sigma$-point-selection core and
state it as a self-contained set-theoretic existence problem. A two-valued state
on $\Lo$ selects, within each block (maximal compatible context, a countable
partition of $\Omega$ into elements of $\Lo$), exactly one ``true'' cell; the
selections across overlapping blocks must agree on shared events; $\sigma$-additivity
constrains the selection over countable joins; non-Dirac means the selection is
realised by no single point.

\begin{question}[$\sigma$-point-selection, bare form]\label{q:bare}
Given a set $\Omega$ and a family $\{\Pi_\alpha\}$ of countable partitions of
$\Omega$ (the blocks of a concrete $\sigma$-OML, pairwise overlapping in their
common refinements), does there exist a selection assigning to each $\Pi_\alpha$
one cell $c_\alpha \in \Pi_\alpha$ such that
\begin{enumerate}[label=\textup{(\alph*)}]
  \item \textup{(coherence)} the selections agree on events common to several
    blocks;
  \item \textup{($\sigma$-additivity)} the induced $\{0,1\}$-valuation commutes
    with countable disjoint unions;
  \item \textup{(non-principality)} there is no $\omega_0 \in \Omega$ with
    $c_\alpha = $ the cell of $\Pi_\alpha$ containing $\omega_0$ for all $\alpha$;
  \item \textup{(prescription)} the selection realises a prescribed finite
    pattern that no single point realises?
\end{enumerate}
\end{question}

This is the object the Localization Theorem isolates: a witness for
Definition~\ref{def:sigma-essential} exists iff Question~\ref{q:bare} has a
positive answer for some admissible carrier. The bare form makes two features
legible that the lattice language obscures.

\begin{remark}[Why this is not simply a measurable cardinal]\label{rem:not-measurable}
If the blocks $\{\Pi_\alpha\}$ were \emph{compatible} --- a single Boolean base ---
a coherent non-principal $\sigma$-additive $\{0,1\}$-selection would be exactly a
non-principal countably-complete ultrafilter, hence a measurable cardinal
\cite{Jech2003}. The defining feature of Question~\ref{q:bare} is that the blocks
are \emph{incompatible}: the coherence condition (a) is imposed across
\emph{overlapping, mutually non-distributive} contexts, and is strictly stronger
than an ultrafilter condition on any single block. Consequently
Question~\ref{q:bare} does not reduce to the existence of a measurable cardinal;
its Boolean shadow does, but the non-distributive coherence requirement is
precisely what the reduction cannot cross. We make this non-reduction precise
against the operator-algebra literature in \S\ref{sec:boundary}.
\end{remark}

\begin{remark}[Relation to sheaf-theoretic contextuality]\label{rem:sheaf}
Question~\ref{q:bare} is the assertion that a system of locally coherent
$\{0,1\}$-valuations admits no \emph{global section} of the corresponding
presheaf --- the sheaf-theoretic signature of contextuality in the sense of
Abramsky and Brandenburger~\cite{AbramskyBrandenburger2011}, in its
$\sigma$-additive form~\cite{AbramskyBarbosaEtAl2022}. The novelty here is not the
sheaf framing, which is standard, but the conjunction of three features that the
existing contextuality literature does not treat together: countable additivity,
a non-Hilbert (abstractly concrete) carrier, and a $\{0,1\}$-valued
(non-probabilistic) selection. It is exactly this conjunction that is open.
\end{remark}


\section{The boundary map}\label{sec:boundary}

We locate Question~\ref{q:bare} by recording where its answer is known and where
the obvious bridges fail.

\paragraph{The Polish boundary: empty above it.}
Derr and Williamson~\cite{DerrWilliamson2023} prove a $\sigma$-extension
criterion for coherent probabilities on pre-Dynkin systems: when the carrier is
\emph{Polish-representable} --- $\Omega$ Polish, the generated $\sigma$-class
equal to the Borel $\sigma$-algebra, with inner-regular blocks --- finite
coherence implies $\sigma$-extendability (their Theorem~D.6, via
Maharam~\cite{Maharam1972}). Specialised to two-valued states, this says that on a
Polish-representable carrier every finitely coherent $s_0$ globalises: \emph{no
$\sigma$-essential contextual state exists}. Hence a witness, if one exists, must
live on a carrier that is \emph{not} Polish-representable. This is the upper
boundary of the problem.

\paragraph{The Hilbert boundary: the analogue exists, but does not transfer.}
Blecher and Weaver~\cite{BlecherWeaver2017} prove that a singular countably
additive \emph{pure} state on the projection lattice of $B(\ell^2(\kappa))$
exists iff $\kappa$ is Ulam measurable --- a genuine large-cardinal-sensitive
state-existence dichotomy on a non-distributive $\sigma$-complete OML. This is the
nearest positive precedent: it shows that countable additivity of states on
non-distributive structures \emph{can} be set-theoretically sensitive. But it does
not settle Question~\ref{q:bare}, for two reasons, neither of which is a reduction
we can complete and both of which point the same way.

First, the carrier is of the wrong kind. The projection lattice of $B(H)$ for
$\dim H \geq 3$ admits \emph{no} two-valued state at all
(Kochen--Specker~\cite{KochenSpecker1967}), so it is not concrete: it is not a
sub-orthoposet of any powerset. A concrete $\sigma$-OML, by definition, has an
order-determining family of two-valued states. The Blecher--Weaver object and ours
are therefore not the same kind of structure, and their dichotomy concerns
\emph{pure} states (the relevant notion on $B(H)$), not two-valued ones, which are
generically absent there by Gleason's theorem~\cite{Gleason1957}.

Second, the obvious transfer mechanism fails. The natural way to carry a
state-existence result between a non-commutative algebra and an abelian (Boolean)
subalgebra is to factor through a maximal abelian subalgebra. For the
\emph{pure}, two-valued case this factorisation does not hold: Akemann and
Weaver~\cite{AkemannWeaver2008} exhibit a pure state on $B(H)$ multiplicative on no
masa, and Blecher--Weaver themselves rely on Marcus--Spielman--Srivastava paving
precisely because the pure state does not reduce to a masa. So one cannot obtain a
positive answer to Question~\ref{q:bare} by importing Blecher--Weaver through a
Boolean subobject.

\begin{remark}\label{rem:no-bridge}
Remarks~\ref{rem:not-measurable} and the two paragraphs above are deliberately
stated at the strength of \emph{cited facts assembled}, not of a non-reducibility
theorem. We do not claim a formal separation in any fixed reduction calculus; we
claim, and the cited results establish, that the two obvious bridges --- ``it is a
measurable cardinal'' and ``it transfers from Blecher--Weaver through a masa'' ---
do not close Question~\ref{q:bare}, and that the obstruction in both cases is the
same non-distributivity. This is what leaves the question open rather than a
corollary of known work.
\end{remark}

\paragraph{Below: faithful representation forces distributivity.}
One might hope to route Question~\ref{q:bare} through a faithful $\sigma$-tribe
representation --- embedding $\Lo$ into a $\sigma$-algebra of honest sets so that
two-valued states become point evaluations. This is unavailable: a faithful
tribe-of-sets representation forces distributivity, hence Booleanness (the
``Floor''; cf.~\cite{PtakPulmannova1991}). The $\sigma$-Loomis--Sikorski theorem
holds for structures with the Riesz Decomposition Property --- $\sigma$-complete
MV-algebras and RDP effect algebras~\cite{Dvurecenskij2000} --- but OMLs lack RDP
($MO_2$ is the standard witness), and the available representation is only a
$\sigma$-epimorphic image that does not separate two-valued points. So the
realization route that would trivialise Question~\ref{q:bare} is blocked, for the
same non-distributive reason.

\begin{remark}[Segregated carriers do not witness]\label{rem:segregated}
A carrier is \emph{segregated} when its contextuality is carried by a central
Boolean factor --- e.g.\ $MO_\omega$ or the product $\prod_n MO_2$, whose
$\sigma$-additive two-valued states all concentrate at points. On such carriers
every finitely coherent $s_0$ is the restriction of a Dirac, so clause (ii) of
Theorem~\ref{thm:localization} fails and there is no witness. A witness must be
\emph{non-segregated}: the incompatibility must be irreducible, not routed through
a central partition. This is why Definition~\ref{def:sigma-essential} requires
irreducibility, and it sharpens the search for a carrier to the non-segregated,
non-Boolean, non-Polish-representable regime.
\end{remark}


\section{The open problem}\label{sec:open}

Assembling \S\ref{sec:localization}--\S\ref{sec:boundary}: the existence of a
$\sigma$-essential contextual state is equivalent to a positive answer to
Question~\ref{q:bare}, on a carrier that is concrete, $\sigma$-complete,
non-Boolean, irreducible, non-segregated, and non-Polish-representable. The answer
is not settled in either direction. It is not a corollary of the Blecher--Weaver
dichotomy (the carrier is the wrong kind and the masa bridge fails), it does not
reduce to a measurable cardinal (the non-distributive coherence condition is
strictly stronger than its Boolean shadow), and it is not trivialised by a
faithful representation (which would force distributivity).

\begin{question}[Main]\label{q:main}
Does there exist, in ZFC or under a large-cardinal hypothesis, a concrete
$\sigma$-complete non-Boolean irreducible OML carrying a $\sigma$-essential
contextual state? Equivalently, is the existence of a non-principal coherent
$\sigma$-additive $\{0,1\}$-selection across overlapping countable partitions
(Question~\ref{q:bare}, non-segregated case) provable, refutable, or independent?
\end{question}

We record what a resolution would require. A positive ZFC answer would be a
concrete construction crossing the non-Polish, non-segregated regime --- which the
present obstructions (Remarks~\ref{rem:not-measurable},~\ref{rem:no-bridge}) show
cannot be assembled from countably many compatible pieces, so it must be globally
non-compositional. A negative answer would extend the Polish-representable
emptiness of Derr--Williamson past its topological hypotheses to all concrete
$\sigma$-OMLs, against the Blecher--Weaver evidence that the analogous Hilbert cell
is inhabited under a large cardinal. An independence result would supply a model
in which the selection exists and one in which it does not; the precedents for
such results in the set theory of measure-theoretic structures (the von
Neumann--Maharam problem~\cite{BalcarJechPazak2005}, real-valued measurable
cardinals~\cite{Fremlin}) are all real-valued and strictly positive, leaving the
two-valued, non-distributive case --- Question~\ref{q:main} --- as the open edge.


\section{A forcing target}\label{sec:forcing}

The shape of the boundary map suggests where a resolution of
Question~\ref{q:main} should be sought: not in ZFC and not in a single named
combinatorial principle, but in an independence proof. We record the target
precisely, as a hand-off for set-theoretic methods.

For an admissible carrier $\Lo$ (concrete, $\sigma$-complete, non-Boolean,
irreducible) write
\[
  \Phi(\Lo) \;:\equiv\; \text{every two-valued state on a finite }\perp\text{-closed }
  B \subseteq \Lo \text{ extends to a global $\sigma$-additive two-valued state,}
\]
the \emph{$\sigma$-additive lifting} property, and let
$\Psi :\equiv \exists\,\Lo\ \neg\Phi(\Lo)$ assert the existence of a witness in the
sense of Definition~\ref{def:sigma-essential}. By Theorem~\ref{thm:localization},
$\neg\Phi(\Lo)$ holds iff some local pattern realises clause~(ii); by
Proposition~\ref{prop:boolean}, $\Phi$ holds for all Boolean $\Lo$, so any witness
is non-distributive. $\Psi$ is the existence sentence whose status is
Question~\ref{q:main}.

Two features make $\Psi$ a forcing target rather than a construction target.
First, by the analysis of \S\ref{sec:bare}--\S\ref{sec:boundary} no admissible
$\Lo$ can be assembled compositionally: any countable, refinable approximation
collapses to the Boolean case (Proposition~\ref{prop:boolean},
Remark~\ref{rem:segregated}), so a witness must be uncountably generated and
globally non-compositional. Second, an independence proof need not exhibit a fixed
$\Lo$: it suffices to produce a model of set theory in which $\Psi$ holds and one
in which it fails, the carrier being supplied by the model rather than constructed
in advance --- as in the von Neumann problem on measure algebras, whose
independence is established by models, not by a fixed algebra
\cite{BalcarJechPazak2005, Maharam1972}.

\begin{remark}[Candidate technology and conjectured strength]\label{rem:technology}
The natural setting is the set theory of measure-theoretic and operator-algebraic
structures. The nearest precedent with an explicit cardinal is Blecher--Weaver's
Ulam-measurable dichotomy on $B(\ell^2(\kappa))$~\cite{BlecherWeaver2017}; we have
argued (\S\ref{sec:boundary}) that it does not transfer, so the strength of $\Psi$
is \emph{conjectural and not inherited}. The relevant cardinal, if any, is
Ulam-measurable rather than real-valued measurable: $\Psi$ concerns two-valued
($\{0,1\}$) selections, where the real-valued-measurable machinery governing
$[0,1]$-valued measures does not directly apply. Boolean-valued models / forcing
over a candidate carrier, and the von Neumann--Maharam circle of independence
results~\cite{BalcarJechPazak2005, Fremlin}, are the natural tools; but all known
precedents are real-valued and strictly positive, and the two-valued
non-distributive case is, as far as we are aware, untreated. We therefore state the
strength as an open question, not a theorem:
\[
  \text{Is } \Psi \text{ independent of ZFC, and if so, of what consistency strength?}
\]
\end{remark}

\paragraph{Acknowledgements.}
The author thanks the anonymous referees (forthcoming) and acknowledges the use of
large language model assistance (Anthropic Claude, OpenAI ChatGPT) for literature
search, verification of cited statements against primary sources, and exposition;
all mathematical claims were independently verified by the author.
